\documentclass[11pt]{article}
\usepackage[margin=1in]{geometry}
\usepackage{amsmath,amssymb,amsthm,mathtools}
\usepackage[T1]{fontenc}
\usepackage{lmodern,microtype,xurl}
\usepackage[colorlinks=true,linkcolor=blue,urlcolor=blue]{hyperref}
\newtheorem{theorem}{Theorem}[section]
\newtheorem{lemma}[theorem]{Lemma}
\newtheorem{proposition}[theorem]{Proposition}
\newtheorem{corollary}[theorem]{Corollary}
\theoremstyle{remark}\newtheorem{remark}[theorem]{Remark}
\newcommand{\E}{\mathbb E}\newcommand{\R}{\mathbb R}
\newcommand{\KL}{\operatorname{D}}\newcommand{\poly}{\operatorname{poly}}
\newcommand{\norm}[1]{\left\lVert#1\right\rVert}
\title{Uniform prior-scale bounds for structured Boolean targets}
\author{ChatGPT-6 Astra \and Kaarel H\"anni \and Dmitry Vaintrob}
\date{25 October 2026}
\begin{document}\maketitle
\begin{abstract}
We assemble target-mass constructions for the original Gaussian tanh
posterior uniformly over all prior prefactors $0<t\le1$, including the
critical endpoint and arbitrarily small prefactors. For each fixed
threshold depth, the sample bound is polynomial jointly in input
dimension, circuit resources, and learner depth at every realizable
width. Fourier variation gives a second polynomial target class at
every learner depth. The proof combines a critical scalar program,
an exact change of Gaussian reference, and a separate short-padding
construction. The finite-width intrinsic learning and critical scalar
inputs are stated explicitly; this note is an integration of audited
research arguments, not a self-contained replacement for those inputs.
\end{abstract}

\section{Model and conclusions}
Let $X_d=\{-1,1\}^d$, let $f:X_d\to\{-1,1\}$, and put $H=L-1\ge1$.
The learner has $H$ hidden tanh layers of common width $N$, augmented
input $(1,x)$, and scalar linear readout. The independent Gaussian
entry variances are
\[
 W^1_{ij}:t^2/d,\qquad W^j_{ik}:t^2/N\ (2\le j\le H),
 \qquad v_i:t^2/N^2,\qquad 0<t\le1.
\]
The posterior conditions only on $|z(x)|>1/2$ for every $x\in X_d$
and on agreement with the noiseless training labels. Write
\[
 A_f=\{w:f(x)z_w(x)>1/2\text{ for all }x\in X_d\},\qquad
 a_{N,t}(f)=-N^{-1}\log\mu_{N,t}(A_f).
\]
A realizable width means $\mu_{N,t}(A_f)>0$. The learning criterion is:
for $0<\varepsilon,\delta<1$, with training probability at least $1-\delta$, the posterior assigns
probability strictly greater than $2/3$ to the correct label outside a
set of test-distribution mass at most $\varepsilon$. The input
distribution is arbitrary.

\begin{theorem}[Threshold circuits, all scales and learner depths]
\label{thm:threshold}
There is a universal constant $C$ with the following property. Suppose
$f$ is computed by a layered threshold circuit with $h\ge1$ gate
stages ending in the output bit $f$, width at most $C_0$, normalized gate margins at least one, and
coefficient magnitudes at most $B$. Constants may be carried as
ordinary circuit gates. Put
\[
 R=d+C_0+B+h+L+4.
\]
For every $L\ge h+1$, every realizable width, and every $0<t\le1$,
the stated learning criterion holds whenever
\[
 n\ge \frac C\varepsilon\left[
 R^{C(h+1)^2}(1+\log(1/\varepsilon))^C+
 \log(1/\delta)\right].                         \tag{1}
\]
If $H-h\ge C[\log(d+2)+h\log R]$, the exponent
$C(h+1)^2$ can be replaced by $C(h+1)$.
\end{theorem}

\begin{theorem}[Fourier variation at every learner depth]
\label{thm:fourier}
Let
\[
 A(f)=\sum_{S\subseteq[d]}|\widehat f(S)|,
 \qquad f=\sum_S\widehat f(S)\prod_{i\in S}x_i.
\]
There are universal $C,c$ such that for every $L\ge2$, every
realizable width, and every $0<t\le1$, it suffices that
\[
 n\ge\frac C\varepsilon\left[
 (d+L+A(f)+2)^c(1+\log(1/\varepsilon))^c+
 \log(1/\delta)\right].                         \tag{2}
\]
In particular, the bound is polynomial for every parity, uniformly in
its degree and in the original Gaussian scale.
\end{theorem}

Neither theorem includes arbitrary polynomial-size circuits of
unrestricted threshold depth with a polynomial bound independent of
$h$. Nor does a tanh teacher become a threshold circuit merely by
replacing its hidden activations by their signs.

\section{Explicit analytic inputs}
Set
\[
 b=t^{-2}-1,\qquad \theta=b/2,\qquad
 c_t=\frac{t^{-2}-1+2\log t}{2}.
\]
Let $V_*$ be the critical scalar Bellman baseline from the critical
depth analysis. Define $\mathcal R(D)\ge1$ by
\[
 \mathcal R(D)-1-\log\mathcal R(D)=2D.
\]
The following two intrinsic learning inputs are proved in the
companion research notes listed at the end.

\begin{proposition}[Uniform intrinsic inputs]
\label{prop:intrinsic}
There are universal polynomial degrees such that the original
posterior sample bound, with $a=a_{N,t}(f)$ or any deterministic
upper bound for it, is polynomial in $d,L,1+t^2a$ and
\[
 P_t(a+t^{-2},H)
 =\max\{1,t^H\mathcal R((a+t^{-2})/H)^{H/2}\},      \tag{3}
\]
apart from the usual factor
$\varepsilon^{-1}(1+\log(1/\varepsilon))^C$ and additive
$C\varepsilon^{-1}\log(1/\delta)$.
There is also a sharper bound, polynomial in $d,L,1+t^2a$ times
$\exp(CJ)$, where
\[
 E=\bigl(a+4/(Ht^2)-Hc_t-V_*+b\bigr)_+,
 \qquad
 J=\min\{\sqrt{HE},E/b\}.                         \tag{4}
\]
At $b=0$, the second argument is omitted. Both statements hold at
every realizable width, uniformly over $0<t\le1$.
\end{proposition}

These inputs use a bounded dual certificate with arbitrarily chosen
action slack, scale-normalized importance quadrature, and a final
split between large widths and the full finite-width VC class. In
particular, they contain no separate inverse power of $t$.

The critical target inputs are as follows. A threshold prefix of
depth $h$ admits deterministic rare groups with seed fraction $p_h$
such that, writing $S=\log(d+2)+h\log R$,
\[
 \log(1/p_h)\le CS,\qquad
 \text{critical prefix cost}\le1/(8L).
                                                               \tag{5}
\]
At every $N\ge\exp(CS)$ its target score is uniformly accurate to
$\exp(-CS)$, after increasing constants, and its actual target
component $p_N$ lies in $[p_h,2p_h]$. Its diagonal is at most one.
For $H-h\ge CS$, the critical continuation can be chosen with charged
scalar costs $D_j$ satisfying
\[
 \text{total action}\le V_*+CS^2/H,\quad
 \sum_j\sqrt{D_j}\le CS,\quad
 \sum_jD_j\le CS,                                \tag{6}
\]
and all tracked components at least $\exp(-CS)$. The final increasing
scalar path, including readout, has cost at most three. The critical
finite-width channel theorem has metric amplification $\exp(CS)$
for this program. These statements are the quantitative finite-prefix,
occupation, and scalar-steering inputs from the critical analysis.

The square-root estimate in (6) concerns the \emph{charged scalar
entropy upper costs}. It cannot be replaced by a smaller vector KL
after data processing when bounding scalar second moments.

\section{Changing scale without changing the physical scalar program}

\begin{lemma}[Exact change of reference]
If $\rho$ is a scalar control law with $m=\E_\rho Z^2$, then
\[
 \KL(\rho\Vert\mathcal N(0,t^2))
 =\KL(\rho\Vert\mathcal N(0,1))+c_t+\theta(m-1).    \tag{7}
\]
Moreover, if $D=\KL(\rho\Vert\mathcal N(0,1))$, then
$(m-1)_+\le C(\sqrt D+D)$.
\end{lemma}
\begin{proof}
Equation (7) is the Gaussian log-density ratio integrated under
$\rho$. The entropy variational principle, applied to the square of
a standard Gaussian, gives $D\ge(m-1-\log m)/2$; inversion for
$m\ge1$ gives the bound.
\end{proof}

\begin{lemma}[Smaller residual noise preserves the target component]
Put $\psi_r(z)=\E\tanh(z+\sqrt rG)$. For $z>0$, $\psi_r(z)$ is
nonnegative and nonincreasing in $r$. A critical Gaussian scalar
stage with the same target component and no larger residual can
therefore be reused at any $t\le1$, while tracking exactly the
critical target component and no larger residual.
\end{lemma}
\begin{proof}
The heat identity gives
$\partial_r\psi_r(z)=\tfrac12\E\tanh''(z+\sqrt rG)\le0$.
Indeed $\tanh''$ is odd and negative on the positive half-line,
and a Gaussian shifted by $z>0$ puts greater density at $u$ than at
$-u$ for $u>0$. Continuity treats $r=0$.

For a Gaussian signal with variance $av$ and residual variance $r$,
the total next diagonal is $F(av+t^2r)$, where
$F(q)=\E\tanh^2(\sqrt qG)$ is increasing. The available conditional
mean component increases when the residual decreases, whereas the
total diagonal decreases. If a normalized score $b_0$ has component
$v_0\ge v$, replace it by
\[
 b=\sqrt{v/v_0}\,b_0+\sqrt{1-v/v_0}\,\xi,
\]
where $\xi$ is an independent uniform sign included under both
reference and controlled laws. This keeps unit score norm and
component exactly $v$, costs zero KL, and increases its pointwise
bound by at most one. The new residual is the diagonal minus $v$.
\end{proof}

\begin{proposition}[Uniform long-padding target mass]
\label{prop:long}
For the threshold target and $S$ above, if $H-h\ge CS$, then for
every $N\ge\exp(CS)$ and every $0<t\le1$,
\[
 a_{N,t}(f)\le Hc_t+V_*+C(\theta S+S^2/H).          \tag{8}
\]
The constants and width threshold are independent of $t$.
\end{proposition}
\begin{proof}
Keep the critical prefix group fractions and mean shifts. Their
Gaussian entropy costs are divided by $t^2$, so the prefix cost is
at most $1/(8Lt^2)$. The Gaussian noise only decreases, preserving
the prefix estimates and allowing the same seed component.

Reuse the physical critical Gaussian scalar laws, with the score
downgrade from the preceding lemma. The final non-Gaussian scalar
path uses the same normalized law as at criticality; its stable-lift
estimate improves when residual noise decreases. Equation (7) and
(6) charge at most one $c_t$ per hidden stage and an additional
$C\theta S$. The readout costs at most three at criticality, so its
factor $t^{-2}=1+2\theta$ adds only $C\theta$. This proves (8) for
the ideal controlled path, with the coefficient of $V_*$ exactly one.

For the finite-width lift, the physical metric multiplier of a
scalar law is at most $\sqrt{\max\{t^2,m\}}\le
\sqrt{\max\{1,m\}}\le\exp(C\sqrt D)$. Thus total amplification
is $\exp(CS)$ independently of $t$. Tracked score bounds and
inverse components are also $\exp(CS)$.
Use action slack $\tau=1/(Ht^2)$ and controlled success loss
\[
 \kappa=\frac{c}{H(1+t^2\overline a)},
\]
where $\overline a$ is the displayed deterministic upper budget.
Since $t^2c_t\le1/2$ and $t^2\theta\le1/2$,
$t^2\overline a\le C(H+S+S^2/H)$. Hence $\kappa^{-1}$ is
polynomial independently of $t$. The channel concentration and
binary-KL comparison therefore require only $N\ge\exp(CS)$.
Their added $1/(Ht^2)=(1+2\theta)/H$ is absorbed in (8).
\end{proof}

Applying (4) to (8) gives $J=O(S)$: if $b\ge S/H$, use $E/b$;
otherwise use $\sqrt{HE}$. All normalized action factors and the
width threshold are at most $\exp(CS)$. This proves the long-padding
part of Theorem~\ref{thm:threshold}.

\section{Short padding}
Let $M=H-h$ and retain the same seed. Fix a universal small
$p_*>0$ and put $S_0=\log(p_*/p_h)$; increase the resource constants
so that $S_0\ge1$ and $S_0\le CS$. If $M=0$, a direct readout has
normalized action at most $C/p_h$.

For $1\le M<CS$, prescribe components growing geometrically from
$p_N$ to $p_*$. A Gaussian scalar variance multiplier in
\[
 1\le a_j\le128\exp(S_0/M)                       \tag{9}
\]
achieves each required increase exactly, uniformly in $t\le1$.
For completeness, if the target component is $v$ and the next one
is $w\le p_*=1/128$, taking $a=128w/v$ gives total preactivation
variance at most two. Gaussian integration by parts and
$\E\operatorname{sech}^2(\sqrt uG)\ge1/6$ for $u\le2$ give
an available component at least $av/36>w$; at $a=1$ it is at most
$v$. Continuity and the zero-cost score downgrade suffice.

The resulting finite-width budget obeys
\[
 t^2\overline a\le C\{1+M\exp(S_0/M)\},\qquad
 \Gamma\le128^{M/2}\exp(S_0/2).                  \tag{10}
\]
For $M=0$, replace the first right side by $C/p_h$.
The same normalized finite-width slack argument gives a threshold
$N\ge\exp(CS)$. Since $\mathcal R(D)\le4D+2$, (3) yields
\[
 \log P_t(\overline a+t^{-2},H)
 \le C(h+1)(S+1).                               \tag{11}
\]
Indeed for $M\ge1$ the upper bound is at most
$C H+C(H/M)S_0$, and $H=h+M$, $M<CS$; the case $M=0$ follows
directly with $H=h$.
Equations (10)--(11) and Proposition~\ref{prop:intrinsic} prove
(1) at widths beyond $\exp(CS)$. All smaller widths are bounded
by that same threshold, so the VC bound for the entire finite-width
tanh class gives (1) there as well. This completes
Theorem~\ref{thm:threshold}.

\section{Fourier seeds and consequences}
The critical Fourier construction supplies, for each parity $\chi_S$,
a mixture of Gaussian mean shifts and cube sign symmetrizations
with mean $c_S\chi_S$, $c_S\ge1/[8(d+1)]$, and entropy at most
\[
 D_*=C(d+1)^3\log^2(d+2).
\]
At variance $t^2/d$, keep the same physical shifts. Their entropy is
at most $D_*/t^2$, the saturation estimate improves, and symmetry
still gives an exact parity mean with the same lower bound. Mixing
the signed parity laws with weights proportional to
$|\widehat f(S)|/c_S(t)$ gives exact mean $f/Z_f(t)$, where
$Z_f(t)\le8(d+1)A(f)$, at the same entropy upper bound.

Mix this controlled law into the raw first law with probability
$\vartheta=(D_*L)^{-1}$. A latent tag score gives the exact component
\[
 v=\frac1{D_*LZ_f(t)^2}\ge
 \frac1{D_*L[8(d+1)A(f)]^2},                     \tag{12}
\]
with prefix cost at most $1/(Lt^2)$ and constant diagonal at most
one. The tag is included under both laws and adds no independent
entropy penalty. Use score downgrade to track the right side of
(12), and initialize the virtual critical scalar comparator with
diagonal one. Bounded scores and features give a polynomial
finite-width seed threshold by concentration over the cube.

Now $h=1$ and $S=O(\log(d+L+A(f)+2))$. The long-padding proof
and the short-padding proof apply unchanged, including $H=1$ with
direct readout. In (11), $h+1$ is an absolute constant. This proves
Theorem~\ref{thm:fourier}.

\begin{corollary}[Two useful regimes]
Polynomial-size Boolean formulas and deterministic or nondeterministic branching programs
of polynomial length and polynomial state width have
$\exp(O(\log^2 L))$ sample bounds when their input dimension and
resources are polynomial in learner depth $L$, uniformly over
$0<t\le1$ and all realizable widths.
\end{corollary}
\begin{proof}
Balanced formulas have logarithmic threshold depth and polynomial
size. For branching programs, compose their Boolean transition
matrices in a balanced tree using
$(AB)_{pq}=\bigvee_r(A_{pr}\wedge B_{rq})$.
This gives threshold depth $O(\log T)$ and size $O(TW^3)$ for
length $T$ and width $W$. Thus $S=O(\log^2 L)$ and the learner
has enough padding for Proposition~\ref{prop:long}.
\end{proof}

\section*{Proof dependencies and audit status}
The local scale-change, residual monotonicity, score downgrade,
short-padding estimates, and their combination with the intrinsic
bounds have been independently checked by the upper-bound,
critical-depth, and complexity-barrier agents. The following scratch
notes identify the imported analytic arguments precisely:
\begin{itemize}
\item \texttt{uniform-scale-intrinsic.md} and
\texttt{uniform-sharp-intrinsic.md}: normalized all-width learning.
\item \path{finite-width-dual-certificate.md} and
\path{uniform-baseline-gain.md}: extraction and baseline gain.
\item \texttt{nearcritical-target-reuse.md}: uniform critical-program reuse.
\item Section 4 of this manuscript: the short-padding schedule and
its explicit gain bound. The critical finite-width channel input below
supplies its concentration and relative-entropy comparison.
\item The critical-depth manuscript's finite circuit prefix,
scalar steering, occupation, finite channel, and controlled Fourier
seed sections: the quantitative critical inputs (5)--(6).
\end{itemize}
The named proof notes accompany this manuscript in the
\path{proof-notes/} directory. Exact links to the critical-depth
source inputs appear in the package README and dependency manifest.
These analytic learning statements are not claimed to be machine verified.
\end{document}
